\subsection{Private Key Encryption}

Recall that for OTPs, $n$ bits of key were required. 
With PRGs, we can create $n$ bits of "good-enough" 
random out of fewer than $n$ bits. The secret seed 
for the PRG takes the place of a key. 

Consider an efficient eavesdropping adversary observing 
the ciphertext. The encryption scheme is secure if the 
adversary can guess better than random with negligible $v(n)$
probability. 

We can define security as having indistinguishable 
encryptions in the presence of an eavesdropper 
(EAV-secure) thus: a symmetric encryption scheme
(Gen, Enc, Dec) is secure if 
$\forall m_0, m_1:\{\text{Enc}_k(m_0)\} \approx_c \{\text{Enc}_k(m_1)\}$
for $k \leftarrow^\mathdollar \{0,1\}^n$. 

Consider: 
\begin{itemize}
    \item $G$ is a PRG
    \item $\text{Gen}(1^n): k\leftarrow^\mathdollar \{0,1\}^n$
    \item $\text{Enc}_k(m) = m \oplus G(k)$
    \item $\text{Dec}_k(c) = c \oplus G(k)$
\end{itemize}
We wish to prove PRG encryption is EAV-secure. 
We may prove by reduction. We want to show 
construction $\Pi$ is secure. We want to show 
that $\Pi$ is secure as long as some underlying 
problem $X$ is hard. The idea is to show that if 
you break $\Pi$ then you also solve $X$. By 
contrapositive, then if $X$ is hard then $\Pi$ 
is also hard. 

Proofs by reduction follow this template:
\begin{itemize}
    \item We want to show a construction $\Pi$ secure 
    \item We want to show that $\Pi$ is secure as long as some underlying problem $X$ is hard
    \item An algorithm that "breaks" $\Pi$ can be turned 
            into an algorithm that "breaks" $X$. 
    \item Therefore, if breaking $\Pi$ would break 
    $X$, then $\Pi$ is secure as long as $X$ is hard. 
    Basically, if breaking $\Pi$ would require an 
    attacker to break $X$, then breaking $\Pi$ is at 
    least as hard as breaking $X$
\end{itemize}

We will study these attack models: 
\begin{itemize}
    \item Ciphertext only attack: eavesdropper observes multiple ciphertexts between Alice and Bob, possibly of the same message, and has prior knowledge of the distribution of messages sent (e.g. like letter frequency in the English language).
    However, ideally an adversary would not be able to tell 
    if the same message is sent. However, if an algorithm 
    is deterministic, then the same message produces 
    the same ciphertext. We will deal with this later
    \item Known plaintext attack: adversary can see plaintext-ciphertext pairs
    \item Chosen plaintext attack: adversary picks what messages it can see the encryptions of
    \item Chosen ciphertext attack: adversary can gather 
    information by obtaining the decryptions of chosen 
    ciphertexts.
\end{itemize}

Model A is stronger than model B, if a cipher secure against 
A is also secure against B. It is a fact that 
CTO < KPA < CPA < CCA. 